% kernel/engineering_objects-DE.tex

\chapter{Vom berechneten Übergang zum Engineering Object}
\label{chap:engineering-objects}

Die Berechnung liefert eine parametrisierte Berlinish-Komposition
$H_{\mathrm{in}}+C+H_{\mathrm{out}}$. Zu diesem Zeitpunkt ist sie ein Vorschlag
im Raum ausdrückbarer Übergangsfunktionen. Die Aussage, dass sie den Übergang
zwischen zwei identifizierten Nachbarelementen verändert, ist von anderer Art:
Der Vorschlag betrifft nun ein bestimmtes Alignment.

\section{Worüber macht der Vorschlag eine Aussage?}

\begin{definition}[Engineering Object]
\label{def:engineering-object}
Ein Engineering Object ist eine dauerhafte ingenieurtechnische Entität, deren
Identität unabhängig von ihren Repräsentationen, Beobachtungen, Realisierungen
und Berechnungsmodellen ist.
\end{definition}

Das Alignment ist der dauerhafte Gegenstand. Der vorgeschlagene Übergang ist
ein Element seiner geordneten Konstruktion und trägt zu einem möglichen Zustand
dieses Alignments bei. Ein SPOT-fähiges Engineering-Object-Modell kann
Gegenstand, Zustand und ihre Beziehungen tragen; ein Softwaredatensatz erzeugt
jedoch nicht ihre Ingenieuridentität. Solver-Ausgabe, daraus erzeugte CAD-Kurve
und Zeile in einer Austauschdatei sind Beschreibungen des Vorschlags, nicht drei
neue Alignments.

Der Vorschlag muss deshalb das betroffene Alignment, den Vorgängerzustand, aus
dem er abgeleitet wurde, sowie Intervall und Nachbarelemente benennen, deren
Konstruktion er ändern würde. Ohne diese Bezüge bleibt das numerische Ergebnis
eine wiederverwendbare Übergangsbeschreibung statt eines instanziierten
Alignment-Elements.

\section{Welche Art von Beschreibung liegt vor?}

Tabelle~\ref{tab:meaning-roles} hält die Unterscheidungen zusammen, die dieser
Teil fortlaufend verwendet. Es sind Beziehungen zu einem gemeinsamen
Ingenieurgegenstand, nicht konkurrierende Namen für dasselbe.

\begin{table}[htbp]
\centering
\small
\caption{Rollen im laufenden Übergangsbeispiel}
\label{tab:meaning-roles}
\begin{tabular}{@{}p{0.15\textwidth}p{0.31\textwidth}p{0.44\textwidth}@{}}
\toprule
Rolle & Im laufenden Beispiel & Was die Rolle nicht begründet \\
\midrule
Objekt & das fortbestehende Alignment & keine bestimmte Geometrie, Datei oder Annahme \\
Repräsentation & CAD-Kurve, abgetastete Polylinie oder Austauschdatensatz & weder Objektidentität noch physische Wahrheit \\
Zustand & beabsichtigter oder durch Krümmungsänderung vorgeschlagener Zustand & keine Autorität allein aufgrund seiner Berechenbarkeit \\
Evidenz & Importgeschichte, Berechnungsprotokoll oder Vermessungsbeobachtung & keine automatische Annahme oder Anwendbarkeit \\
Bewertung & Prüfung gegen Geschwindigkeit, Überhöhung, Realisierung und Regeln & keine Ingenieurentscheidung \\
Autorität & verantwortliche Annahme, Ablehnung oder Auswahl & keine Übertragung der Verantwortung an das Werkzeug \\
\bottomrule
\end{tabular}
\end{table}

Damit kann aus einem importierten Übergang eine abgeleitete editierbare
Repräsentation entstehen, ohne den Import zu zerstören. Der importierte
Datensatz bleibt quellengebundene Evidenz; Rekonstruktion erzeugt einen
nachvollziehbar abgeleiteten Zustand; Bearbeitung erzeugt einen Vorschlag. Jede
Beschreibung behält ihre Provenienz und ihren Wahrheitsmodus und bezieht sich,
wenn die Identitätsaussage begründet ist, auf dasselbe Alignment.

\section{Warum als Nächstes Identität erforderlich ist}

Die Objektunterscheidung verhindert, dass Datei, Kurve oder Solver-Ergebnis zum
Alignment werden. Sie sagt noch nicht, ob der Austausch des Klothoidenkerns,
die Änderung eines Halbwellparameters oder der Neuaufbau einer importierten
Kurve dasselbe Objekt bewahrt. Diese Frage verlangt Identität durch Änderung.
